% =====================================================================
%  Anatomía de un LLM — Fase III — Módulo 15
%  Construye tu GPT
%  Compilar: latexmk -r ../latexmkrc M15-mini-gpt.tex
% =====================================================================
\documentclass[11pt,a4paper,twoside,openany]{book}
\usepackage{../estilo-apuntes}

\newcommand{\ApuntesCurso}{Anatomía de un LLM}
\newcommand{\ApuntesModulo}{Módulo 15}

% Notación propia del módulo
\newcommand{\tok}[1]{\texttt{#1}}
\newcommand{\vs}{\texttt{\textvisiblespace}}
\newcommand{\nl}{\texttt{\textbackslash n}}
\newcommand{\LN}{\operatorname{LN}}
\newcommand{\MHA}{\operatorname{MHA}}
\newcommand{\MLP}{\operatorname{MLP}}
\newcommand{\GELU}{\operatorname{GELU}}
\newcommand{\clip}{\operatorname{clip}}
\newcommand{\KL}{D_{\mathrm{KL}}}
\newcommand{\Var}{\operatorname{Var}}
\newcommand{\Cov}{\operatorname{Cov}}
\newcommand{\topk}{\operatorname{top}_k}
\newcommand{\ind}{\mathbb{1}}
\newcommand{\Qm}{\bm{Q}}
\newcommand{\Km}{\bm{K}}
\newcommand{\Vm}{\bm{V}}
\newcommand{\Mm}{\bm{M}}
\makeatletter\g@addto@macro\UrlBreaks{\do\-}\makeatother

\hypersetup{pdftitle={Módulo 15. Construye tu GPT (resumen)},pdfauthor={Anatomía de un LLM}}

\makeatletter
\newlabel{ej:corpus}{{1.2}{1}{El corpus de la práctica}{ejemplo.2}{}}
\newlabel{ej:corpus@cref}{{[ejemplo][2][]1.2}{[1][1][]1}}
\newlabel{ej:sal-y-sol}{{1.5}{2}{}{ejemplo.5}{}}
\newlabel{ej:sal-y-sol@cref}{{[ejemplo][5][]1.5}{[1][2][]2}}
\newlabel{ej:gradiente}{{1.11}{3}{}{ejemplo.11}{}}
\newlabel{ej:gradiente@cref}{{[ejemplo][11][]1.11}{[1][3][]3}}
\newlabel{eje:1-desplazada}{{1.1}{4}{}{ejercicio.1.1}{}}
\newlabel{eje:1-desplazada@cref}{{[ejercicio][1][1]1.1}{[1][4][]4}}
\newlabel{eje:1-hola}{{1.2}{4}{}{ejercicio.1.2}{}}
\newlabel{eje:1-hola@cref}{{[ejercicio][2][1]1.2}{[1][4][]4}}
\newlabel{eje:1-d2l}{{1.3}{4}{}{ejercicio.1.3}{}}
\newlabel{eje:1-d2l@cref}{{[ejercicio][3][1]1.3}{[1][4][]4}}
\newlabel{eje:1-estable}{{1.4}{5}{}{ejercicio.1.4}{}}
\newlabel{eje:1-estable@cref}{{[ejercicio][4][1]1.4}{[1][4][]5}}
\newlabel{eje:1-ppl}{{1.5}{5}{}{ejercicio.1.5}{}}
\newlabel{eje:1-ppl@cref}{{[ejercicio][5][1]1.5}{[1][5][]5}}
\newlabel{eje:1-cabeza-cero}{{1.6}{5}{}{ejercicio.1.6}{}}
\newlabel{eje:1-cabeza-cero@cref}{{[ejercicio][6][1]1.6}{[1][5][]5}}
\newlabel{fig:minigpt}{{2.1}{7}{Flujo de datos del mini-GPT de la \cref {def:minigpt}. Las flechas de la izquierda son las conexiones residuales; cada subcapa normaliza su entrada (pre-norm)}{figure.2.1}{}}
\newlabel{fig:minigpt@cref}{{[figure][1][2]2.1}{[1][7][]7}}
\newlabel{ej:mascara}{{2.4}{8}{}{ejemplo.4}{}}
\newlabel{ej:mascara@cref}{{[ejemplo][4][]2.4}{[1][8][]8}}
\newlabel{eje:2-sin-mascara}{{2.1}{10}{}{ejercicio.2.1}{}}
\newlabel{eje:2-sin-mascara@cref}{{[ejercicio][1][2]2.1}{[1][10][]10}}
\newlabel{eje:2-mascara}{{2.2}{10}{}{ejercicio.2.2}{}}
\newlabel{eje:2-mascara@cref}{{[ejercicio][2][2]2.2}{[1][10][]10}}
\newlabel{eje:2-param32}{{2.3}{10}{}{ejercicio.2.3}{}}
\newlabel{eje:2-param32@cref}{{[ejercicio][3][2]2.3}{[1][10][]10}}
\newlabel{eje:2-cabezas}{{2.4}{10}{}{ejercicio.2.4}{}}
\newlabel{eje:2-cabezas@cref}{{[ejercicio][4][2]2.4}{[1][10][]10}}
\newlabel{eje:2-gpt2}{{2.5}{11}{}{ejercicio.2.5}{}}
\newlabel{eje:2-gpt2@cref}{{[ejercicio][5][2]2.5}{[1][10][]11}}
\newlabel{eje:2-permutacion}{{2.6}{11}{}{ejercicio.2.6}{}}
\newlabel{eje:2-permutacion@cref}{{[ejercicio][6][2]2.6}{[1][11][]11}}
\newlabel{ej:encoge}{{3.4}{13}{}{ejemplo.4}{}}
\newlabel{ej:encoge@cref}{{[ejemplo][4][]3.4}{[1][13][]13}}
\newlabel{ej:calendario}{{3.7}{14}{}{ejemplo.7}{}}
\newlabel{ej:calendario@cref}{{[ejemplo][7][]3.7}{[1][14][]14}}
\newlabel{ej:coste}{{3.11}{15}{}{ejemplo.11}{}}
\newlabel{ej:coste@cref}{{[ejemplo][11][]3.11}{[1][15][]15}}
\newlabel{eje:3-orden}{{3.1}{16}{}{ejercicio.3.1}{}}
\newlabel{eje:3-orden@cref}{{[ejercicio][1][3]3.1}{[1][15][]16}}
\newlabel{eje:3-calendario}{{3.2}{16}{}{ejercicio.3.2}{}}
\newlabel{eje:3-calendario@cref}{{[ejercicio][2][3]3.2}{[1][16][]16}}
\newlabel{eje:3-cs336-cal}{{3.3}{16}{}{ejercicio.3.3}{}}
\newlabel{eje:3-cs336-cal@cref}{{[ejercicio][3][3]3.3}{[1][16][]16}}
\newlabel{eje:3-clip}{{3.4}{16}{}{ejercicio.3.4}{}}
\newlabel{eje:3-clip@cref}{{[ejercicio][4][3]3.4}{[1][16][]16}}
\newlabel{eje:3-sgd}{{3.5}{16}{}{ejercicio.3.5}{}}
\newlabel{eje:3-sgd@cref}{{[ejercicio][5][3]3.5}{[1][16][]16}}
\newlabel{eje:3-flops}{{3.6}{16}{}{ejercicio.3.6}{}}
\newlabel{eje:3-flops@cref}{{[ejercicio][6][3]3.6}{[1][16][]16}}
\newlabel{eje:3-proyeccion}{{3.7}{16}{}{ejercicio.3.7}{}}
\newlabel{eje:3-proyeccion@cref}{{[ejercicio][7][3]3.7}{[1][16][]16}}
\newlabel{ej:entropia-corpus}{{4.7}{19}{Entropía del corpus sintético}{ejemplo.7}{}}
\newlabel{ej:entropia-corpus@cref}{{[ejemplo][7][]4.7}{[1][19][]19}}
\newlabel{eje:4-estado}{{4.1}{20}{}{ejercicio.4.1}{}}
\newlabel{eje:4-estado@cref}{{[ejercicio][1][4]4.1}{[1][20][]20}}
\newlabel{eje:4-fijos}{{4.2}{20}{}{ejercicio.4.2}{}}
\newlabel{eje:4-fijos@cref}{{[ejercicio][2][4]4.2}{[1][20][]20}}
\newlabel{eje:4-linea}{{4.3}{20}{}{ejercicio.4.3}{}}
\newlabel{eje:4-linea@cref}{{[ejercicio][3][4]4.3}{[1][20][]20}}
\newlabel{eje:4-brecha}{{4.4}{20}{}{ejercicio.4.4}{}}
\newlabel{eje:4-brecha@cref}{{[ejercicio][4][4]4.4}{[1][20][]20}}
\newlabel{eje:4-memoria}{{4.5}{21}{}{ejercicio.4.5}{}}
\newlabel{eje:4-memoria@cref}{{[ejercicio][5][4]4.5}{[1][21][]21}}
\newlabel{eje:4-condicional}{{4.6}{21}{}{ejercicio.4.6}{}}
\newlabel{eje:4-condicional@cref}{{[ejercicio][6][4]4.6}{[1][21][]21}}
\newlabel{ej:temperatura}{{5.5}{23}{}{ejemplo.5}{}}
\newlabel{ej:temperatura@cref}{{[ejemplo][5][]5.5}{[1][23][]23}}
\newlabel{eje:5-conceptos}{{5.1}{26}{}{ejercicio.5.1}{}}
\newlabel{eje:5-conceptos@cref}{{[ejercicio][1][5]5.1}{[1][25][]26}}
\newlabel{eje:5-topk}{{5.2}{26}{}{ejercicio.5.2}{}}
\newlabel{eje:5-topk@cref}{{[ejercicio][2][5]5.2}{[1][26][]26}}
\newlabel{eje:5-temp}{{5.3}{26}{}{ejercicio.5.3}{}}
\newlabel{eje:5-temp@cref}{{[ejercicio][3][5]5.3}{[1][26][]26}}
\newlabel{eje:5-topp}{{5.4}{26}{}{ejercicio.5.4}{}}
\newlabel{eje:5-topp@cref}{{[ejercicio][4][5]5.4}{[1][26][]26}}
\newlabel{eje:5-conmuta}{{5.5}{26}{}{ejercicio.5.5}{}}
\newlabel{eje:5-conmuta@cref}{{[ejercicio][5][5]5.5}{[1][26][]26}}
\newlabel{eje:5-contexto}{{5.6}{26}{}{ejercicio.5.6}{}}
\newlabel{eje:5-contexto@cref}{{[ejercicio][6][5]5.6}{[1][26][]26}}
\newlabel{ap:soluciones}{{A}{27}{Soluciones de los ejercicios}{appendix.A}{}}
\newlabel{ap:soluciones@cref}{{[appendix][1][2147483647]A}{[1][27][]27}}
\newlabel{sol:eje:1-desplazada}{{A.1}{27}{Tema 1}{section*.6}{}}
\newlabel{sol:eje:1-desplazada@cref}{{[subappendix][1][2147483647,1]A.1}{[1][27][]27}}
\newlabel{sol:eje:1-hola}{{A.1}{27}{Tema 1}{section*.7}{}}
\newlabel{sol:eje:1-hola@cref}{{[subappendix][1][2147483647,1]A.1}{[1][27][]27}}
\newlabel{sol:eje:1-d2l}{{A.1}{28}{Tema 1}{section*.8}{}}
\newlabel{sol:eje:1-d2l@cref}{{[subappendix][1][2147483647,1]A.1}{[1][28][]28}}
\newlabel{sol:eje:1-estable}{{A.1}{28}{Tema 1}{section*.9}{}}
\newlabel{sol:eje:1-estable@cref}{{[subappendix][1][2147483647,1]A.1}{[1][28][]28}}
\newlabel{sol:eje:1-ppl}{{A.1}{29}{Tema 1}{section*.10}{}}
\newlabel{sol:eje:1-ppl@cref}{{[subappendix][1][2147483647,1]A.1}{[1][29][]29}}
\newlabel{sol:eje:1-cabeza-cero}{{A.1}{29}{Tema 1}{section*.11}{}}
\newlabel{sol:eje:1-cabeza-cero@cref}{{[subappendix][1][2147483647,1]A.1}{[1][29][]29}}
\newlabel{sol:eje:2-sin-mascara}{{A.2}{30}{Tema 2}{section*.12}{}}
\newlabel{sol:eje:2-sin-mascara@cref}{{[subappendix][2][2147483647,1]A.2}{[1][29][]30}}
\newlabel{sol:eje:2-mascara}{{A.2}{30}{Tema 2}{section*.13}{}}
\newlabel{sol:eje:2-mascara@cref}{{[subappendix][2][2147483647,1]A.2}{[1][30][]30}}
\newlabel{sol:eje:2-param32}{{A.2}{31}{Tema 2}{section*.14}{}}
\newlabel{sol:eje:2-param32@cref}{{[subappendix][2][2147483647,1]A.2}{[1][31][]31}}
\newlabel{sol:eje:2-cabezas}{{A.2}{31}{Tema 2}{section*.15}{}}
\newlabel{sol:eje:2-cabezas@cref}{{[subappendix][2][2147483647,1]A.2}{[1][31][]31}}
\newlabel{sol:eje:2-gpt2}{{A.2}{31}{Tema 2}{section*.16}{}}
\newlabel{sol:eje:2-gpt2@cref}{{[subappendix][2][2147483647,1]A.2}{[1][31][]31}}
\newlabel{sol:eje:2-permutacion}{{A.2}{32}{Tema 2}{section*.17}{}}
\newlabel{sol:eje:2-permutacion@cref}{{[subappendix][2][2147483647,1]A.2}{[1][32][]32}}
\newlabel{sol:eje:3-orden}{{A.3}{33}{Tema 3}{section*.18}{}}
\newlabel{sol:eje:3-orden@cref}{{[subappendix][3][2147483647,1]A.3}{[1][33][]33}}
\newlabel{sol:eje:3-calendario}{{A.3}{33}{Tema 3}{section*.19}{}}
\newlabel{sol:eje:3-calendario@cref}{{[subappendix][3][2147483647,1]A.3}{[1][33][]33}}
\newlabel{sol:eje:3-cs336-cal}{{A.3}{34}{Tema 3}{section*.20}{}}
\newlabel{sol:eje:3-cs336-cal@cref}{{[subappendix][3][2147483647,1]A.3}{[1][34][]34}}
\newlabel{sol:eje:3-clip}{{A.3}{34}{Tema 3}{section*.21}{}}
\newlabel{sol:eje:3-clip@cref}{{[subappendix][3][2147483647,1]A.3}{[1][34][]34}}
\newlabel{sol:eje:3-sgd}{{A.3}{35}{Tema 3}{section*.22}{}}
\newlabel{sol:eje:3-sgd@cref}{{[subappendix][3][2147483647,1]A.3}{[1][35][]35}}
\newlabel{sol:eje:3-flops}{{A.3}{35}{Tema 3}{section*.23}{}}
\newlabel{sol:eje:3-flops@cref}{{[subappendix][3][2147483647,1]A.3}{[1][35][]35}}
\newlabel{sol:eje:3-proyeccion}{{A.3}{35}{Tema 3}{section*.24}{}}
\newlabel{sol:eje:3-proyeccion@cref}{{[subappendix][3][2147483647,1]A.3}{[1][35][]35}}
\newlabel{sol:eje:4-estado}{{A.4}{36}{Tema 4}{section*.25}{}}
\newlabel{sol:eje:4-estado@cref}{{[subappendix][4][2147483647,1]A.4}{[1][36][]36}}
\newlabel{sol:eje:4-fijos}{{A.4}{36}{Tema 4}{section*.26}{}}
\newlabel{sol:eje:4-fijos@cref}{{[subappendix][4][2147483647,1]A.4}{[1][36][]36}}
\newlabel{sol:eje:4-linea}{{A.4}{37}{Tema 4}{section*.27}{}}
\newlabel{sol:eje:4-linea@cref}{{[subappendix][4][2147483647,1]A.4}{[1][37][]37}}
\newlabel{sol:eje:4-brecha}{{A.4}{37}{Tema 4}{section*.28}{}}
\newlabel{sol:eje:4-brecha@cref}{{[subappendix][4][2147483647,1]A.4}{[1][37][]37}}
\newlabel{sol:eje:4-memoria}{{A.4}{37}{Tema 4}{section*.29}{}}
\newlabel{sol:eje:4-memoria@cref}{{[subappendix][4][2147483647,1]A.4}{[1][37][]37}}
\newlabel{sol:eje:4-condicional}{{A.4}{38}{Tema 4}{section*.30}{}}
\newlabel{sol:eje:4-condicional@cref}{{[subappendix][4][2147483647,1]A.4}{[1][38][]38}}
\newlabel{sol:eje:5-conceptos}{{A.5}{38}{Tema 5}{section*.31}{}}
\newlabel{sol:eje:5-conceptos@cref}{{[subappendix][5][2147483647,1]A.5}{[1][38][]38}}
\newlabel{sol:eje:5-topk}{{A.5}{39}{Tema 5}{section*.32}{}}
\newlabel{sol:eje:5-topk@cref}{{[subappendix][5][2147483647,1]A.5}{[1][39][]39}}
\newlabel{sol:eje:5-temp}{{A.5}{39}{Tema 5}{section*.33}{}}
\newlabel{sol:eje:5-temp@cref}{{[subappendix][5][2147483647,1]A.5}{[1][39][]39}}
\newlabel{sol:eje:5-topp}{{A.5}{40}{Tema 5}{section*.34}{}}
\newlabel{sol:eje:5-topp@cref}{{[subappendix][5][2147483647,1]A.5}{[1][40][]40}}
\newlabel{sol:eje:5-conmuta}{{A.5}{40}{Tema 5}{section*.35}{}}
\newlabel{sol:eje:5-conmuta@cref}{{[subappendix][5][2147483647,1]A.5}{[1][40][]40}}
\newlabel{sol:eje:5-contexto}{{A.5}{41}{Tema 5}{section*.36}{}}
\newlabel{sol:eje:5-contexto@cref}{{[subappendix][5][2147483647,1]A.5}{[1][41][]41}}
\makeatother

\begin{document}
\frontmatter
\begin{titlepage}
\centering
\vspace*{2.5cm}
{\large\scshape Anatomía de un LLM}\par
\vspace{0.4cm}
{\normalsize Fundamentos matemáticos de los modelos de lenguaje}\par
\vspace{2.2cm}
\rule{\textwidth}{0.8pt}\par\vspace{0.5cm}
{\small\scshape Fase III \textperiodcentered\ Secuencias y Transformer}\par\vspace{0.3cm}
{\Huge\bfseries Módulo 15\par}\vspace{0.3cm}
{\LARGE Construye tu GPT\par}
\vspace{0.2cm}
{\large Entrenamiento, evaluación y generación de un Transformer decoder\par}
\vspace{0.5cm}\rule{\textwidth}{0.4pt}\par
\vspace{1.6cm}
{\large Resumen del módulo}\par
\vfill
{\small Edición de octubre de 2026}\par
\end{titlepage}

\chapter*{Sobre este resumen}
\addcontentsline{toc}{chapter}{Sobre este resumen}
\markboth{Sobre este resumen}{}

Este documento reúne lo indispensable de los apuntes del Módulo 15: el texto expositivo, todas las definiciones y todos los resultados (teoremas, proposiciones, lemas y corolarios), las tablas y las observaciones de contenido. Se han suprimido las demostraciones, los ejemplos, las lecturas, las figuras, los ejercicios y sus soluciones. La numeración de temas, secciones, enunciados y ecuaciones es la misma que en los apuntes, de modo que cualquier referencia puede buscarse allí; las que remiten a elementos suprimidos (un ejemplo, una figura, una ecuación de una demostración) conservan su número en los apuntes.

\section*{Notación}
\begin{center}\small
\begin{tabularx}{\textwidth}{@{}l>{\raggedright\arraybackslash}X@{}}
\toprule
Símbolo & Significado \\
\midrule
$\V$, $K=\abs{\V}$ & vocabulario y su tamaño \\
$u_{1:n}$ & corpus tokenizado de $n$ tokens \\
$T$, $B$ & longitud de la ventana de contexto y tamaño del lote \\
$x_{1:T}$, $y_{1:T}$ & ventana de entrada y ventana objetivo, con $y_t=x_{t+1}$ \\
$d$, $L$, $n_h$, $d_k=d/n_h$ & dimensión del modelo, número de bloques, de cabezas y dimensión por cabeza \\
$\Qm,\Km,\Vm$, $\Mm$ & matrices de consultas, claves y valores, y máscara causal \\
$z\in\R^{K}$, $p=\softmax(z)$ & logits y distribución predicha en una posición \\
$\ell$, $\mathcal{L}$ & pérdida de una posición y pérdida media \\
$\theta$, $N$ & parámetros del modelo y su número \\
$s$, $S$, $W$ & paso de entrenamiento, número total de pasos y pasos de calentamiento \\
$\eta_s$, $\eta_{\max}$, $\eta_{\min}$ & tasa de aprendizaje en el paso $s$ y sus valores máximo y mínimo \\
$\lambda$, $\beta_1,\beta_2$ & decaimiento de pesos y coeficientes de los momentos de AdamW \\
$c$, $\clip_c$ & umbral y operador de recorte del gradiente \\
$\tau$, $k$, $\rho$ & temperatura, parámetro de top-$k$ y umbral de top-$p$ \\
$H(p)$, $H(p,q)$, $\KL(p\Vert q)$ & entropía, entropía cruzada y divergencia de Kullback--Leibler (en nats) \\
\bottomrule
\end{tabularx}
\end{center}


\tableofcontents

\mainmatter
% =====================================================================
\chapter{Datos, ventanas y objetivo de entrenamiento}
\label{tema:1}
% =====================================================================

Un modelo de lenguaje autorregresivo se entrena con una única secuencia larga de enteros, de la que se extraen ventanas de longitud fija. Este tema formaliza esa preparación de los datos, define la función que se minimiza (la entropía cruzada media sobre todas las posiciones de un lote) y demuestra tres propiedades de esa función que se usan durante todo el módulo: su gradiente respecto de los logits, su invariancia frente a traslaciones de los logits, que permite calcularla de forma estable, y su valor cuando la predicción es uniforme. Se sigue el §4.1 y el §5.1 de la tarea~1 de CS336~\cite{cs336}, el §9.3.3 de~\cite{d2l} y la clase~7 de CS224N~\cite{cs224n}.

\section{El corpus como secuencia de tokens}

\begin{definicion}[Corpus tokenizado]\label{def:corpus}
Sea $\V$ un vocabulario de tamaño $K$, identificado con $\{0,1,\dots,K-1\}$. Un \emph{corpus tokenizado} es una secuencia $u_{1:n}=(u_1,\dots,u_n)\in\V^{n}$. Si el texto original está formado por varios documentos, se concatenan en una sola secuencia, separados por un token delimitador.
\end{definicion}

La concatenación es una decisión de ingeniería con una justificación precisa: todas las ventanas tienen la misma longitud y no hace falta rellenar con símbolos de relleno (\emph{padding}), que consumen cómputo sin aportar señal~\cite[§5.1]{cs336}. El precio es que una ventana puede cruzar la frontera entre dos documentos; el delimitador permite al modelo aprender que lo que sigue no depende de lo anterior.

\addtocounter{definicion}{1}%

Un corpus sintético tiene una ventaja metodológica: se conoce el proceso que lo genera. Se sabe, por ejemplo, que tras «receta.» viene uno de seis verbos y que «un reloj» exige «negro» y no «negra». Esto permite juzgar qué ha aprendido el modelo y calcular la pérdida mínima alcanzable (\cref{tema:4}).

\section{Ventanas de contexto y lotes}

\begin{definicion}[Ventana de contexto]\label{def:ventana}
Sea $T\ge 1$ la \emph{longitud de contexto}. Para cada $i\in\{1,\dots,n-T\}$, la \emph{ventana} que empieza en $i$ es el par
\[
x=(u_i,u_{i+1},\dots,u_{i+T-1}),\qquad y=(u_{i+1},u_{i+2},\dots,u_{i+T}),
\]
es decir, $y_t=x_{t+1}$ para $t<T$ e $y_T=u_{i+T}$. La secuencia $y$ es la \emph{ventana objetivo}.
\end{definicion}

\begin{proposicion}\label{prop:ventanas}
Un corpus de $n$ tokens tiene exactamente $n-T$ ventanas de longitud $T$, y cada una contiene $T$ problemas de predicción: predecir $y_t$ a partir del prefijo $x_{1:t}$, para $t=1,\dots,T$.
\end{proposicion}

\addtocounter{definicion}{1}%

\begin{definicion}[Lote]\label{def:lote}
Un \emph{lote} de tamaño $B$ es un par de matrices de enteros $X,Y\in\V^{B\times T}$ cuyas filas son $B$ ventanas, con inicios elegidos de forma independiente y uniforme en $\{1,\dots,n-T\}$. Un lote contiene $B\cdot T$ predicciones.
\end{definicion}

Con $B=T=32$, cada paso de entrenamiento ve $1024$ predicciones; en $800$ pasos son $819\,200$, unas $57$ veces los $14\,459$ tokens de entrenamiento del \cref{ej:corpus}. Con tan pocos datos es inevitable repetir; los grandes modelos, en cambio, ven cada token una vez o pocas veces (módulo~16).

\begin{observacion}[Partición en entrenamiento y validación]\label{obs:particion}
El conjunto de validación se obtiene cortando el corpus \emph{por posición}: el 90~\% inicial ($14\,459$ tokens) se usa para entrenar y el 10~\% final ($1\,607$ tokens) para validar. Si se eligieran ventanas al azar para cada conjunto, una ventana de validación podría solaparse con una de entrenamiento y la pérdida de validación subestimaría el error de generalización.
\end{observacion}

La alternativa al muestreo independiente de inicios es partir el corpus en trozos consecutivos sin solape, con un desplazamiento inicial aleatorio en cada época, como hace~\cite[§9.3.3]{d2l}. El \cref{eje:1-d2l} compara ambas estrategias.

\section{La pérdida de entropía cruzada}

En cada posición $t$ de cada ventana, el modelo devuelve un vector de \emph{logits} $z_t\in\R^{K}$ y la distribución predicha $p_t=\softmax(z_t)$, con $p_{t,j}=e^{z_{t,j}}/\sum_{a}e^{z_{t,a}}$.

\begin{definicion}[Pérdida de entrenamiento]\label{def:perdida}
La pérdida de una posición con objetivo $y\in\V$ y logits $z$ es
\[
\ell(z,y)=-\ln\softmax(z)_y.
\]
La pérdida de un lote es la media sobre sus $B\cdot T$ posiciones,
\[
\mathcal{L}(\theta)=\frac{1}{BT}\sum_{b=1}^{B}\sum_{t=1}^{T}\ell\bigl(z_{b,t}(\theta),\,Y_{b,t}\bigr),
\]
donde $z_{b,t}(\theta)$ son los logits que el modelo de parámetros $\theta$ produce en la posición $t$ de la fila $b$ a partir de $X_{b,1:t}$.
\end{definicion}

La pérdida $\ell$ es la entropía cruzada entre la distribución degenerada en $y$ y la distribución predicha~\cite[§4.1]{cs336}. Su media sobre el corpus es la estimación de máxima verosimilitud del módulo~1 aplicada a un modelo con parámetros compartidos: minimizar $\mathcal{L}$ equivale a maximizar $\sum\ln P_\theta(y_t\mid x_{1:t})$.

\begin{definicion}[Perplejidad]\label{def:ppl}
Si una secuencia de $m$ predicciones tiene pérdidas $\ell_1,\dots,\ell_m$, su \emph{perplejidad} es
\[
\PPL=\exp\Bigl(\frac1m\sum_{i=1}^{m}\ell_i\Bigr).
\]
\end{definicion}

La perplejidad es el número de opciones equiprobables entre las que el modelo duda en media: un modelo uniforme sobre $K$ símbolos tiene perplejidad $K$. Es independiente de la base del logaritmo, porque $e^{\mathcal{L}}=2^{\mathcal{L}/\ln 2}$. En la práctica del módulo, el modelo de bigramas del módulo~1 llega a $2{,}035$ nats por carácter ($\PPL=7{,}65$) y el mini-GPT a $0{,}324$ ($\PPL=1{,}38$).

\section{Gradiente y estabilidad numérica}

\begin{proposicion}[Gradiente de la entropía cruzada]\label{prop:gradiente}
Sea $p=\softmax(z)$. Para todo $i\in\V$,
\[
\frac{\partial\ell(z,y)}{\partial z_i}=p_i-\ind[i=y],
\]
es decir, $\nabla_z\ell=p-e_y$, donde $e_y$ es el vector \emph{one-hot} de $y$.
\end{proposicion}

\addtocounter{definicion}{1}%

El gradiente de la \cref{prop:gradiente} es la única señal que recibe todo el mini-GPT: la retropropagación (módulo~10) la transporta desde los logits hasta los embeddings.

\begin{proposicion}[Invariancia frente a traslaciones]\label{prop:traslacion}
Para todo $c\in\R$, $\softmax(z+c\mathbf 1)=\softmax(z)$. En consecuencia, con $z^{\star}=\max_a z_a$,
\[
\ell(z,y)=-(z_y-z^{\star})+\ln\sum_{a}e^{\,z_a-z^{\star}},
\]
y todos los exponentes del segundo miembro son menores o iguales que~0.
\end{proposicion}

La fórmula estable evita el desbordamiento: en precisión simple el mayor número representable es unos $3{,}4\cdot 10^{38}\approx e^{88{,}7}$, de modo que $e^{1000}$ daría infinito y el cociente, un valor indefinido. Es la técnica \emph{log-sum-exp} del módulo~11, que CS336 exige en su implementación de la entropía cruzada~\cite[§4.1]{cs336}.

\begin{proposicion}[Pérdida de la predicción uniforme]\label{prop:lnK}
Si todos los logits de una posición son iguales, la pérdida vale $\ln K$ para cualquier objetivo.
\end{proposicion}

Con $K=46$, $\ln 46=3{,}829$ nats ($5{,}52$ bits). Una red recién inicializada con pesos pequeños produce logits casi iguales, y su pérdida inicial está cerca de ese valor; en la práctica del módulo el primer lote da $4{,}046$. Una pérdida inicial muy superior a $\ln K$ indica una inicialización defectuosa: el modelo está muy seguro de respuestas equivocadas.


% =====================================================================
\chapter{La arquitectura mini-GPT}
\label{tema:2}
% =====================================================================

El mini-GPT es el bloque Transformer del módulo~14 apilado $L$ veces entre una capa de entrada, que convierte tokens y posiciones en vectores, y una capa de salida, que devuelve un logit por token del vocabulario. Ninguna pieza es nueva; este tema fija la definición completa del modelo, demuestra la propiedad de causalidad que justifica entrenar las $T$ predicciones de una ventana con una sola pasada y cuenta sus parámetros. La arquitectura es la de GPT-2~\cite[§2.3]{gpt2}, construida sobre el decoder de Vaswani \emph{et al.}~\cite{vaswani}; se sigue también el §3 de la tarea~1 de CS336~\cite{cs336}, la clase~5 de CS224N~\cite{cs224n} y el §11.9.3 de~\cite{d2l}.

\section{Definición del modelo}

Se recuerdan del módulo~11 la \emph{normalización de capa} de un vector $v\in\R^{d}$,
\[
\LN(v)=\gamma\odot\frac{v-\mu(v)\mathbf 1}{\sqrt{\sigma^{2}(v)+\varepsilon}}+\beta,\qquad \mu(v)=\frac1d\sum_{i}v_i,\quad \sigma^{2}(v)=\frac1d\sum_i\bigl(v_i-\mu(v)\bigr)^{2},
\]
con parámetros $\gamma,\beta\in\R^{d}$, y del módulo~14 el perceptrón por posición $\MLP(v)=W_2\,\GELU(W_1v+b_1)+b_2$, con $W_1\in\R^{4d\times d}$ y $W_2\in\R^{d\times 4d}$.

\begin{definicion}[Atención causal multicabeza]\label{def:mha}
Sea $H\in\R^{T\times d}$ la matriz cuyas filas son las representaciones de las $T$ posiciones, y sea $n_h$ un divisor de $d$, con $d_k=d/n_h$. Se calculan $\Qm=HW_Q$, $\Km=HW_K$ y $\Vm=HW_V$, con $W_Q,W_K,W_V\in\R^{d\times d}$, y se parten por columnas en $n_h$ bloques $\Qm_r,\Km_r,\Vm_r\in\R^{T\times d_k}$. La \emph{máscara causal} es la matriz $\Mm\in\{0,-\infty\}^{T\times T}$ con $\Mm_{tj}=0$ si $j\le t$ y $\Mm_{tj}=-\infty$ si $j>t$. Cada cabeza calcula
\[
\mathrm{head}_r=\softmax\Bigl(\frac{\Qm_r\Km_r^{\top}}{\sqrt{d_k}}+\Mm\Bigr)\Vm_r,\qquad r=1,\dots,n_h,
\]
con la softmax aplicada por filas y el convenio $e^{-\infty}=0$, y la salida es
$\MHA(H)=[\mathrm{head}_1\;\cdots\;\mathrm{head}_{n_h}]\,W_O+\mathbf 1 b_O^{\top}$, con $W_O\in\R^{d\times d}$ y $b_O\in\R^{d}$.
\end{definicion}

En la implementación, las tres proyecciones se agrupan en una sola matriz $d\times 3d$ sin sesgo, como sugiere CS336~\cite[§3.4]{cs336}; el número de parámetros es el mismo.

\begin{definicion}[Mini-GPT]\label{def:minigpt}
Sean $K$, $T$, $d$, $L$ y $n_h$ enteros positivos con $n_h\mid d$. El \emph{mini-GPT} con estos hiperparámetros transforma una ventana $x_{1:T}\in\V^{T}$ en logits $z_{1:T}\in\R^{T\times K}$ mediante:
\begin{enumerate}[label=(\arabic*)]
  \item \emph{entrada}: $h^{0}_t=E_{x_t}+P_t$, donde $E\in\R^{K\times d}$ es la tabla de embeddings de token y $P\in\R^{T\times d}$ la de embeddings de posición, ambas aprendidas;
  \item \emph{bloques}: para $\ell=1,\dots,L$, con $H^{\ell-1}$ la matriz de filas $h^{\ell-1}_t$,
  \[
  A^{\ell}=H^{\ell-1}+\MHA_{\ell}\bigl(\LN_{\ell,1}(H^{\ell-1})\bigr),\qquad
  H^{\ell}=A^{\ell}+\MLP_{\ell}\bigl(\LN_{\ell,2}(A^{\ell})\bigr),
  \]
  donde $\LN$ y $\MLP$ se aplican a cada fila por separado;
  \item \emph{salida}: $z_t=W_{\mathrm{s}}\,\LN_{\mathrm{f}}(h^{L}_t)+b_{\mathrm{s}}$, con $W_{\mathrm{s}}\in\R^{K\times d}$ y $b_{\mathrm{s}}\in\R^{K}$.
\end{enumerate}
La distribución predicha en la posición $t$ es $p_t=\softmax(z_t)$.
\end{definicion}

La forma $x\leftarrow x+f(\LN(x))$ se llama \emph{pre-norm}: la normalización se aplica a la entrada de cada subcapa y el flujo residual queda sin normalizar, lo que estabiliza el entrenamiento de redes profundas (módulo~11). GPT-2 introdujo exactamente estos dos cambios respecto de GPT: la normalización en la entrada de cada subbloque y una normalización final tras el último bloque~\cite[§2.3]{gpt2}. La \cref{fig:minigpt} resume el flujo de datos. Suele decirse que la atención es la etapa en la que las posiciones se \emph{comunican} y el MLP aquella en la que cada posición \emph{procesa} lo recibido.

\addtocounter{figure}{1}%

\section{La máscara causal}

\begin{lema}\label{lem:mascara}
Sea $S\in\R^{T\times T}$ una matriz de puntuaciones. La fila $t$ de $\softmax(S+\Mm)$ vale
\[
a_{tj}=\begin{cases}\dfrac{e^{S_{tj}}}{\sum_{j'\le t}e^{S_{tj'}}} & j\le t,\\[3mm] 0 & j>t,\end{cases}
\]
es decir, es la softmax de $(S_{t1},\dots,S_{tt})$ completada con ceros. En particular, cada fila suma~1 y la fila~1 es $(1,0,\dots,0)$.
\end{lema}

\addtocounter{definicion}{1}%

\begin{teorema}[Causalidad]\label{teo:causal}
En el mini-GPT, para cada $t\in\{1,\dots,T\}$, los logits $z_t$ dependen solo de $x_{1:t}$: si dos ventanas $x$ y $x'$ coinciden en las $t$ primeras posiciones, entonces $z_t(x)=z_t(x')$.
\end{teorema}

\begin{corolario}\label{cor:una-pasada}
Una sola evaluación del mini-GPT sobre la ventana $x_{1:T}$ produce, en cada posición $t$, los mismos logits que se obtendrían evaluándolo sobre el prefijo $x_{1:t}$. Por tanto, una pasada calcula las $T$ distribuciones $P_\theta(\cdot\mid x_{1:t})$ de la \cref{def:perdida}, cada una con su contexto correcto.
\end{corolario}

\begin{observacion}[Fuga del futuro]\label{obs:fuga}
Sin la máscara, la posición $t$ puede atender a $x_{t+1}=y_t$. Una cabeza que copie el token siguiente lleva la pérdida a casi cero sin aprender a predecir nada, y el modelo falla al generar, cuando el futuro no existe. Una pérdida de validación próxima a cero en los primeros cientos de pasos es el síntoma característico (\cref{tab:sintomas}). La implementación original pone a $-\infty$ las entradas de la softmax correspondientes a conexiones ilegales~\cite[§3.2.3]{vaswani}.
\end{observacion}

\section{Recuento de parámetros}

\begin{proposicion}[Parámetros del mini-GPT]\label{prop:parametros}
Con proyecciones $W_Q,W_K,W_V$ sin sesgo, sesgos en $W_O$, en el MLP y en la cabeza, y sin compartir $E$ con $W_{\mathrm{s}}$, el mini-GPT de la \cref{def:minigpt} tiene
\[
N=\underbrace{Kd}_{E}+\underbrace{Td}_{P}+L\,(12d^{2}+10d)+\underbrace{2d}_{\LN_{\mathrm{f}}}+\underbrace{dK+K}_{\text{cabeza}}
\]
parámetros. El número de cabezas $n_h$ no interviene.
\end{proposicion}

\begin{table}[h]
\centering\small
\caption{Parámetros del mini-GPT de la práctica: $K=46$, $T=32$, $d=64$, $L=2$, $n_h=4$.}\label{tab:parametros}
\begin{tabular}{@{}llr@{}}
\toprule
Pieza & Fórmula & Parámetros \\
\midrule
Embedding de token & $Kd$ & 2\,944 \\
Embedding de posición & $Td$ & 2\,048 \\
Atención (por bloque) & $3d^{2}+d^{2}+d$ & 16\,448 \\
MLP (por bloque) & $8d^{2}+5d$ & 33\,088 \\
Dos normalizaciones (por bloque) & $4d$ & 256 \\
Bloque completo & $12d^{2}+10d$ & 49\,792 \\
Dos bloques & $2\,(12d^{2}+10d)$ & 99\,584 \\
Normalización final & $2d$ & 128 \\
Cabeza & $dK+K$ & 2\,990 \\
\midrule
Total & & 107\,694 \\
\bottomrule
\end{tabular}
\end{table}

El 92~\% de los parámetros está en los bloques ($99\,584/107\,694=0{,}925$). El término $12d^{2}$ crece con el cuadrado de la anchura, mientras que las tablas crecen linealmente en $d$; en los modelos grandes la proporción de los bloques es aún mayor, salvo que el vocabulario sea muy grande (\cref{eje:2-gpt2}). Es la misma cuenta que el módulo~14 hace para un modelo de $7\cdot 10^{9}$ parámetros.

\begin{observacion}[Pesos atados]\label{obs:atados}
Vaswani \emph{et al.} comparten la matriz de embeddings con la transformación lineal previa a la softmax~\cite[§3.4]{vaswani}, de modo que $W_{\mathrm{s}}=E$ y se ahorran $Kd$ parámetros. El mini-GPT de la práctica no los ata; con atadura tendría $107\,694-2\,944=104\,750$ parámetros.
\end{observacion}

\section{Del mini-GPT a los modelos abiertos}

Los modelos abiertos más citados usan la misma estructura con algunas piezas sustituidas. LLaMA normaliza la entrada de cada subcapa con RMSNorm, sustituye la GELU por SwiGLU con dimensión oculta $\tfrac23\cdot 4d$ y elimina los embeddings de posición absolutos en favor de las rotaciones RoPE en cada capa~\cite[§2.2]{llama}; CS336 adopta esas mismas piezas~\cite[§3]{cs336}. GPT-2, además de la normalización en la entrada de cada subbloque, escala por $1/\sqrt{N_{\mathrm{r}}}$ al inicializar los pesos de las capas residuales, con $N_{\mathrm{r}}$ el número de capas residuales, usa un vocabulario de $50\,257$ tokens y un contexto de $1024$~\cite[§2.3]{gpt2}. El \cref{tab:arquitecturas} resume las diferencias.

\begin{table}[h]
\centering\small
\caption{Componentes del mini-GPT, de GPT-2 y de LLaMA, según los artículos citados y la implementación de GPT-2 de nanoGPT~\cite{nanogpt}.}\label{tab:arquitecturas}
\begin{tabularx}{\textwidth}{@{}l>{\raggedright\arraybackslash}X>{\raggedright\arraybackslash}X>{\raggedright\arraybackslash}X@{}}
\toprule
Componente & Mini-GPT & GPT-2~\cite{gpt2} & LLaMA~\cite{llama} \\
\midrule
Posiciones & tabla aprendida $T\times d$ & tabla aprendida & RoPE en cada capa \\
Normalización & LayerNorm pre-norm & LayerNorm pre-norm y final & RMSNorm pre-norm \\
MLP & GELU, $d\to 4d\to d$ & GELU, $d\to 4d\to d$ & SwiGLU, dimensión $\tfrac23\cdot 4d$ \\
Tokenizador & caracteres, $K=46$ & BPE de bytes, $K=50\,257$ & BPE (SentencePiece) \\
\bottomrule
\end{tabularx}
\end{table}

Sobre los modelos comerciales, la información pública es limitada: el informe técnico de Gemini indica que sus modelos se construyen sobre decoders Transformer~\cite[§2]{gemini}; xAI publicó en 2024 los pesos y el código de Grok-1, un Transformer de 64 capas con mezcla de 8 expertos y RoPE~\cite{grok1}, y Anthropic no ha publicado la arquitectura de Claude.


% =====================================================================
\chapter{Optimización: AdamW, calendario y recorte}
\label{tema:3}
% =====================================================================

Entrenar el mini-GPT consiste en repetir un paso muy breve: tomar un lote, calcular la pérdida, retropropagar y actualizar los parámetros. Lo que separa un entrenamiento estable de uno que diverge son tres piezas que rodean ese paso: el optimizador AdamW, el calendario de la tasa de aprendizaje y el recorte del gradiente. Este tema las define tal como las especifica la tarea~1 de CS336~\cite[§4.2--4.5]{cs336}, demuestra sus propiedades principales y estima el coste del entrenamiento. AdamW y el calendario coseno se estudiaron en el módulo~7; aquí se usan en su forma de implementación, siguiendo a Loshchilov y Hutter~\cite{adamw,sgdr} y la receta documentada por LLaMA~\cite[§2.3]{llama}.

\section{El paso de entrenamiento}

\begin{definicion}[Paso de entrenamiento]\label{def:paso}
Dados los parámetros $\theta_{s}$ y el estado del optimizador en el paso $s=0,1,\dots,S-1$, un \emph{paso de entrenamiento} consiste en:
\begin{enumerate}[label=\arabic*.]
  \item fijar la tasa $\eta_s$ según el calendario;
  \item extraer un lote $(X,Y)$ (\cref{def:lote});
  \item calcular los logits y la pérdida $\mathcal{L}(\theta_s)$ (\cref{def:perdida});
  \item poner a cero los gradientes acumulados y calcular $g=\nabla_\theta\mathcal{L}(\theta_s)$ por retropropagación;
  \item recortar $g$ por norma global (\cref{def:recorte});
  \item aplicar un paso de AdamW con tasa $\eta_s$ (\cref{def:adamw-impl});
  \item cada cierto número de pasos, evaluar la pérdida de entrenamiento y de validación y guardar un punto de control si la validación mejora (\cref{tema:4}).
\end{enumerate}
\end{definicion}

El orden importa. El recorte actúa sobre el gradiente ya calculado y antes de que el optimizador lo incorpore a sus momentos; si se aplicara después, un gradiente anómalo ya habría contaminado $m$ y $v$. En PyTorch, el paso~4 son las llamadas \texttt{zero\_grad} y \texttt{backward}, y el paso~6 es \texttt{step}~\cite[§4.2.1]{cs336}.

\section{AdamW}

\begin{definicion}[AdamW, forma de implementación]\label{def:adamw-impl}
Sean $\beta_1,\beta_2\in[0,1)$, $\lambda\ge 0$ y $\varepsilon>0$, y sean $m_0=v_0=0$. En la iteración $t=1,2,\dots$, con tasa $\alpha$ y gradiente $g_t$, AdamW realiza
\begin{align*}
\alpha_t&=\alpha\,\frac{\sqrt{1-\beta_2^{t}}}{1-\beta_1^{t}}, &
\theta&\leftarrow\theta-\alpha\lambda\theta,\\
m_t&=\beta_1 m_{t-1}+(1-\beta_1)g_t, &
v_t&=\beta_2 v_{t-1}+(1-\beta_2)g_t^{2},\\
\theta&\leftarrow\theta-\alpha_t\,\frac{m_t}{\sqrt{v_t}+\varepsilon}, &&
\end{align*}
donde el cuadrado, la raíz y el cociente se aplican componente a componente.
\end{definicion}

Es el algoritmo~1 de CS336~\cite[§4.3]{cs336}, que reproduce el algoritmo~2 de Loshchilov y Hutter~\cite{adamw}. La corrección de sesgo de los momentos del módulo~7 aparece aquí agrupada en $\alpha_t$; la diferencia con la forma $\hat m_t/(\sqrt{\hat v_t}+\varepsilon)$ es solo la escala efectiva de $\varepsilon$. El término $-\alpha\lambda\theta$ es el \emph{decaimiento de pesos desacoplado}: no pasa por los momentos y encoge todos los pesos en la misma proporción.

\begin{proposicion}[Encogimiento por decaimiento]\label{prop:encoge}
Si un parámetro recibe gradiente nulo durante los pasos $s=0,\dots,S-1$, con tasas $\eta_s$ y $\lambda\eta_s<1$, su valor final es
\[
\theta_S=\theta_0\prod_{s=0}^{S-1}(1-\lambda\eta_s)\;\le\;\theta_0\exp\Bigl(-\lambda\sum_{s=0}^{S-1}\eta_s\Bigr)\qquad(\theta_0\ge 0).
\]
\end{proposicion}

\addtocounter{definicion}{1}%

Los momentos $m$ y $v$ forman el \emph{estado} del optimizador: dos números por parámetro, además del contador $t$. Por eso el estado ocupa el doble que el modelo y debe guardarse en los puntos de control (\cref{tema:4}). LLaMA usa $\beta_1=0{,}9$, $\beta_2=0{,}95$ y $\lambda=0{,}1$~\cite[§2.3]{llama}; CS336 señala que $(0{,}9;\allowbreak\,0{,}999)$ es la elección típica y $(0{,}9;\allowbreak\,0{,}95)$ la habitual en modelos de lenguaje grandes~\cite[§4.3]{cs336}.

\section{Calentamiento lineal y decaimiento coseno}

\begin{definicion}[Calendario del módulo]\label{def:calendario}
Sean $0\le\eta_{\min}<\eta_{\max}$, $W\ge 1$ pasos de calentamiento y $S>W$ pasos totales. Para $s=0,\dots,S-1$,
\[
\eta_s=\begin{cases}
\eta_{\max}\,\dfrac{s+1}{W}, & s<W,\\[3mm]
\eta_{\min}+\tfrac12(\eta_{\max}-\eta_{\min})\Bigl(1+\cos\dfrac{\pi(s-W)}{S-W}\Bigr), & s\ge W.
\end{cases}
\]
\end{definicion}

El tramo coseno es la ecuación~(5) de SGDR~\cite[§3]{sgdr} con un único ciclo. La práctica usa $\eta_{\min}=\eta_{\max}/10$, como LLaMA~\cite[§2.3]{llama} y la configuración de ejemplo de nanoGPT~\cite{nanogpt}.

\begin{proposicion}[Propiedades del calendario]\label{prop:calendario}
En el calendario de la \cref{def:calendario}:
\begin{enumerate}
  \item $\eta_s$ crece estrictamente en $\{0,\dots,W-1\}$, con $\eta_0=\eta_{\max}/W$ y $\eta_{W-1}=\eta_W=\eta_{\max}$;
  \item $\eta_s$ decrece estrictamente en $\{W,\dots,S-1\}$, y la expresión del tramo coseno vale $\eta_{\min}$ en $s=S$;
  \item en el punto medio del tramo coseno, $s=W+(S-W)/2$, vale $(\eta_{\max}+\eta_{\min})/2$;
  \item para todo $u\in[0,S-W]$, $\eta(W+u)+\eta(S-u)=\eta_{\max}+\eta_{\min}$, donde $\eta(\cdot)$ denota la expresión del tramo coseno evaluada en un real.
\end{enumerate}
\end{proposicion}

\addtocounter{definicion}{1}%

\begin{observacion}[El convenio de CS336]\label{obs:cs336-cal}
CS336 define el calentamiento como $\alpha_t=(t/T_w)\,\alpha_{\max}$ para $t<T_w$, el coseno entre $T_w$ y $T_c$ y una tasa constante $\alpha_{\min}$ para $t>T_c$~\cite[§4.4]{cs336}. Con ese convenio, $\alpha_0=0$ (el primer paso no mueve los pesos, salvo por el decaimiento, que también se multiplica por $\alpha$) y el pico se alcanza en $t=T_w$; con el del módulo, la tasa empieza en $\eta_{\max}/W$ y llega al pico en $s=W-1$. Ambas formas difieren en un paso y son equivalentes en la práctica.
\end{observacion}

La justificación de ambas fases se dio en el módulo~7. Durante el calentamiento, los momentos de Adam se han estimado con muy pocos gradientes, y estos son grandes e irregulares; una tasa grande desde el primer paso puede llevar a una región de la que el entrenamiento no se recupera. Al final, una tasa pequeña reduce el exceso de pérdida debido al ruido del gradiente estocástico. LLaMA documenta 2000 pasos de calentamiento y decaimiento coseno hasta el 10~\% de la tasa máxima~\cite[§2.3]{llama}.

\section{Recorte del gradiente por norma global}

\begin{definicion}[Recorte por norma global]\label{def:recorte}
Sea $g=(g^{(1)},\dots,g^{(P)})$ el gradiente de todos los tensores de parámetros, considerado como un único vector, y $\norm{g}_2=\bigl(\sum_{q}\norm{g^{(q)}}_2^{2}\bigr)^{1/2}$ su \emph{norma global}. Dados $c>0$ y $\varepsilon>0$ pequeño,
\[
\clip_c(g)=\begin{cases} g, & \norm{g}_2<c,\\[1mm] \dfrac{c}{\norm{g}_2+\varepsilon}\,g, & \norm{g}_2\ge c.\end{cases}
\]
\end{definicion}

Es la definición de CS336, con $\varepsilon=10^{-6}$~\cite[§4.5]{cs336}; LLaMA usa $c=1$~\cite[§2.3]{llama}, como la práctica del módulo. El recorte protege de los gradientes anómalos que aparecen cuando la superficie de pérdida tiene acantilados~\cite[§8.2.4]{gbc}, y que de otro modo dañarían durante muchos pasos las medias móviles de Adam.

\begin{proposicion}[Propiedades del recorte]\label{prop:recorte}
Para todo $g$:
\begin{enumerate}
  \item $\clip_c(g)$ es un múltiplo positivo de $g$, luego tiene su misma dirección y sentido;
  \item $\norm{\clip_c(g)}_2<c$ si $\norm{g}_2\ge c$, y $\clip_c(g)=g$ en otro caso;
  \item recortar cada tensor por separado con el mismo umbral produce, en general, un vector que no es múltiplo de $g$.
\end{enumerate}
\end{proposicion}

El recorte global no cambia la dirección de descenso; solo limita la longitud del paso. El recorte por tensor, en cambio, altera el peso relativo de cada parte del modelo en la actualización.

\section{El coste del entrenamiento}

Por la regla del módulo~14, entrenar un modelo de $N$ parámetros con $D$ tokens cuesta $C\approx 6ND$ operaciones de coma flotante (FLOP): $2ND$ en la pasada hacia delante y el doble en la retropropagación. CS336 usa la misma hipótesis sobre la retropropagación~\cite[§4.3, problema \emph{adamw\_accounting}]{cs336}.

\addtocounter{definicion}{1}%


% =====================================================================
\chapter{Evaluación, puntos de control y curvas de pérdida}
\label{tema:4}
% =====================================================================

La pérdida de entrenamiento de un lote es una cantidad ruidosa que, además, no mide lo que interesa: cómo se comporta el modelo con texto que no ha visto. Este tema define el estimador de la pérdida que se usa para dibujar las curvas, describe qué debe contener un punto de control para poder reanudar el entrenamiento sin perturbarlo, demuestra que la entropía del proceso que genera el texto es una cota inferior de la pérdida esperada y presenta una guía de diagnóstico a partir de la forma de las curvas. Se sigue el §7.8 y el capítulo~11 de~\cite{gbc}, los §5.2 y §7.2.3 de la tarea~1 de CS336~\cite{cs336} y el capítulo~2 de Cover y Thomas~\cite{ct}.

\section{Estimar la pérdida}

\begin{definicion}[Pérdida estimada]\label{def:estimada}
Sean $(X^{(1)},Y^{(1)}),\dots,(X^{(r)},Y^{(r)})$ lotes extraídos de un conjunto de datos (entrenamiento o validación). La \emph{pérdida estimada} del modelo de parámetros $\theta$ es
\[
\hat{\mathcal{L}}(\theta)=\frac1r\sum_{q=1}^{r}\mathcal{L}_q(\theta),
\]
donde $\mathcal{L}_q$ es la pérdida media del lote $q$ (\cref{def:perdida}). En la práctica del módulo, $r=20$ y los lotes se extraen con un generador de semilla fija, de modo que son los mismos en todas las evaluaciones.
\end{definicion}

Si los lotes se extraen de forma independiente, $\hat{\mathcal{L}}(\theta)$ es un estimador insesgado de la pérdida esperada y su varianza es la de un lote dividida por $r$. Usar siempre los mismos lotes tiene una ventaja adicional cuando se comparan dos puntos de la curva.

\begin{proposicion}[Lotes comunes]\label{prop:lotes-comunes}
Sean $\theta_a$ y $\theta_b$ dos estados del modelo, y sean $\hat{\mathcal{L}}_a$ y $\hat{\mathcal{L}}_b$ sus pérdidas estimadas, vistas como variables aleatorias a través de la elección de los lotes. Entonces
\[
\Var\bigl(\hat{\mathcal{L}}_a-\hat{\mathcal{L}}_b\bigr)=\Var(\hat{\mathcal{L}}_a)+\Var(\hat{\mathcal{L}}_b)-2\Cov(\hat{\mathcal{L}}_a,\hat{\mathcal{L}}_b).
\]
Si los dos estimadores usan lotes independientes, la covarianza es nula; si usan los mismos lotes y la covarianza es positiva, la varianza de la diferencia es menor.
\end{proposicion}

Dos estados próximos del modelo encuentran difíciles los mismos lotes, de modo que sus pérdidas sobre lotes comunes están positivamente correlacionadas; la diferencia entre dos puntos de la curva refleja entonces sobre todo el cambio del modelo y no el azar del muestreo. Es la técnica de los \emph{números aleatorios comunes} de la simulación estocástica.

Durante la evaluación se usan dos instrucciones de PyTorch: \texttt{model.eval()}, que desactiva los componentes que se comportan distinto al entrenar (como el \emph{dropout}), y el contexto \texttt{torch.no\_grad()}, que evita construir el grafo de cómputo y ahorra memoria y tiempo. Al terminar, \texttt{model.train()} devuelve el modelo al modo de entrenamiento.

\section{Puntos de control}

\begin{definicion}[Punto de control]\label{def:checkpoint}
Un \emph{punto de control} es un registro que contiene todo lo necesario para reanudar el entrenamiento como si no se hubiera interrumpido: los parámetros $\theta$, el estado del optimizador (en AdamW, los momentos $m$, $v$ y el contador de pasos) y el número de paso $s$, que fija la tasa del calendario~\cite[§5.2]{cs336}. Es habitual añadir la pérdida de validación y la configuración del modelo.
\end{definicion}

Si se reanuda sin el estado del optimizador, AdamW reinicia $m=v=0$ y $t=0$: el primer paso usa momentos estimados con un solo gradiente, sin la protección de un calentamiento, y la curva de pérdida da un salto. Si se reanuda sin el número de paso, el calendario vuelve a empezar.

\begin{proposicion}[Memoria del estado de entrenamiento]\label{prop:memoria}
En precisión simple (4~bytes por número), los parámetros, los gradientes y los dos momentos de AdamW de un modelo de $N$ parámetros ocupan $16N$ bytes, sin contar las activaciones.
\end{proposicion}

Para el mini-GPT son $107\,694\cdot 16=1{,}72$~MB, de los que el modelo ocupa $0{,}43$~MB. Para un modelo de $7\cdot 10^{9}$ parámetros, la misma cuenta da $112$~GB, más que la memoria de una GPU de $80$~GB; por eso existen las técnicas de partición del estado del optimizador del módulo~17.

\begin{observacion}[Parada temprana]\label{obs:parada}
Guardar el punto de control solo cuando la pérdida de validación mejora y devolver al final el mejor es el algoritmo de \emph{parada temprana} de Goodfellow, Bengio y Courville~\cite[§7.8]{gbc}, que actúa como regularizador: limita el número efectivo de pasos de optimización. Los grandes entrenamientos guardan además puntos de control frecuentes porque, con miles de GPU, las interrupciones son rutinarias: el informe de Llama~3 documenta 466 interrupciones en un periodo de 54 días de preentrenamiento~\cite[§3.3.4]{llama3}.
\end{observacion}

\section{La cota inferior de la pérdida}

\begin{teorema}[La entropía como cota inferior]\label{teo:cota}
Sean $p$ y $q$ distribuciones sobre $\V$ con $q(x)>0$ siempre que $p(x)>0$. Entonces
\[
H(p,q)=-\sum_{x}p(x)\ln q(x)=H(p)+\KL(p\Vert q)\;\ge\;H(p),
\]
con igualdad si y solo si $q=p$.
\end{teorema}

Si $p$ es la distribución real del token siguiente y $q$ la del modelo, la pérdida esperada no puede bajar de $H(p)$. El \cref{eje:4-condicional} extiende el resultado a distribuciones condicionadas al contexto, que es el caso de un modelo de lenguaje. En consecuencia, si la pérdida media de validación sale claramente por debajo de la entropía del proceso generador, hay una fuga de información. En texto humano no se conoce $H$, pero el mismo resultado explica por qué las curvas de los grandes modelos se aplanan: las leyes de escalado del módulo~16 incluyen un término constante que se interpreta como la entropía irreducible del texto.

\addtocounter{definicion}{1}%

\section{Lectura de las curvas de pérdida}

Una curva de pérdida se interpreta respecto de tres referencias horizontales: $\ln K$, la pérdida de la predicción uniforme (\cref{prop:lnK}); la pérdida del modelo de bigramas, que es lo mínimo que hay que superar; y, cuando se conoce, la entropía $H$ del texto, que no puede superarse (\cref{teo:cota}). El \cref{tab:curva} recoge la curva de la práctica.

\begin{table}[h]
\centering\small
\caption{Pérdidas del mini-GPT base de la práctica, en nats por carácter. Referencias: $\ln 46=3{,}829$; bigramas, $2{,}035$; entropía del corpus, $\approx 0{,}125$.}\label{tab:curva}
\begin{tabularx}{\textwidth}{@{}lcc>{\raggedright\arraybackslash}X@{}}
\toprule
Momento & Entrenamiento & Validación & Lectura \\
\midrule
Paso 0 & $4{,}046$ (un lote) & --- & cerca de $\ln K$: predicción casi uniforme \\
Paso 100 & $2{,}017$ & $2{,}046$ & alcanza el nivel del bigrama \\
Paso 200 & $1{,}040$ & $1{,}154$ & aprende palabras y plantillas \\
Paso 400 & $0{,}370$ & $0{,}425$ & usa el contexto largo \\
Paso 800 & $0{,}264$ & $0{,}324$ & afinado final con tasa baja \\
\bottomrule
\end{tabularx}
\end{table}

\paragraph{La forma.} Hay una caída rápida hasta el nivel del bigrama (pasos 0--100), en la que el modelo aprende las frecuencias de pares de caracteres, y una segunda bajada larga (100--500), en la que aprende a usar el contexto: qué plantilla está escribiendo y qué palabras han aparecido ya. La pendiente final se suaviza a medida que el coseno reduce la tasa.

\paragraph{La brecha.} Al final, la diferencia entre validación y entrenamiento es $0{,}06$ nats. Es un sobreajuste leve: el modelo ha memorizado algunas frases, lo que es esperable cuando cada token se ve unas 57 veces. Si la pérdida de validación empezara a subir mientras la de entrenamiento baja, habría que parar, usar más datos o regularizar más.

\begin{table}[h]
\centering\small
\caption{Síntomas frecuentes en las curvas de pérdida y su diagnóstico probable.}\label{tab:sintomas}
\begin{tabularx}{\textwidth}{@{}>{\raggedright\arraybackslash}X>{\raggedright\arraybackslash}X>{\raggedright\arraybackslash}X@{}}
\toprule
Síntoma & Diagnóstico probable & Qué probar \\
\midrule
No baja de $\ln K$ & error de programación: objetivos desalineados, tasa nula & sobreajustar un solo lote \\
Baja y después da NaN & tasa demasiado alta & bajar $\eta_{\max}$, alargar el calentamiento, recortar \\
Baja muy despacio & tasa demasiado baja & subir $\eta_{\max}$ \\
Validación sube y entrenamiento baja & sobreajuste & más datos, decaimiento de pesos, parar antes \\
Pérdida casi nula en ambas & fuga del futuro & revisar la máscara causal (\cref{obs:fuga}) \\
\bottomrule
\end{tabularx}
\end{table}

La primera prueba de la tercera columna es la estrategia de depuración más recomendada: un modelo correcto debe poder llevar a casi cero la pérdida de un conjunto de datos diminuto, como un único lote; si no lo consigue, el fallo está en el software y no en la capacidad del modelo~\cite[§11.5]{gbc}, \cite[§7.2.3]{cs336}. La misma sección de~\cite{gbc} recomienda razonar con los errores de entrenamiento y de prueba: error bajo en entrenamiento y alto en prueba indica sobreajuste; error alto en ambos deja abiertas la falta de capacidad y un defecto del programa.


% =====================================================================
\chapter{Generación e hiperparámetros}
\label{tema:5}
% =====================================================================

Generar texto es aplicar la regla de la cadena del módulo~1 con el modelo como calculadora de probabilidades condicionadas: se predice la distribución del token siguiente, se muestrea, se añade el token al contexto y se repite. Este tema formaliza el bucle de generación, estudia las tres transformaciones de la distribución que exponen las interfaces de los modelos actuales (temperatura, top-$k$ y top-$p$) y demuestra que la entropía de la distribución crece con la temperatura. Termina con el efecto de los hiperparámetros sobre el entrenamiento del mini-GPT. Se sigue el §6 de la tarea~1 de CS336~\cite{cs336}, el §9.5.5 de~\cite{d2l}, los artículos de Fan, Lewis y Dauphin~\cite{fan} y de Holtzman \emph{et al.}~\cite{holtzman}, y el §11.4 de~\cite{gbc}.

\section{Generación autorregresiva}

\begin{definicion}[Bucle de generación]\label{def:generacion}
Sea un modelo con contexto máximo $T$, una secuencia inicial (\emph{prompt}) $w_{1:m}$, un número máximo de tokens nuevos $G$ y una transformación $\Phi$ que convierte unos logits en una distribución. Para $g=1,\dots,G$:
\begin{enumerate}[label=\arabic*.]
  \item recortar el contexto a sus últimos $T$ tokens, $c=w_{\max(1,\,m+g-T):\,m+g-1}$;
  \item evaluar el modelo sobre $c$ y tomar los logits $z$ de la última posición;
  \item muestrear $w_{m+g}\sim\Phi(z)$ y añadirlo a la secuencia;
  \item detenerse si $w_{m+g}$ es el token delimitador.
\end{enumerate}
\end{definicion}

Por el \cref{cor:una-pasada}, los logits de la última posición son los de la distribución $P_\theta(\cdot\mid c)$. Con $\Phi=\softmax$ se muestrea del propio modelo~\cite[§6, ecuaciones~(21)--(22)]{cs336}; si en cada paso se elige el token más probable, el procedimiento es la \emph{decodificación voraz}~\cite[§9.5.5]{d2l}. Cada paso exige una pasada completa del modelo sobre el contexto; el módulo~19 estudia cómo evitar recalcular lo ya calculado con la caché de claves y valores.

\section{Temperatura}

\begin{definicion}[Softmax con temperatura]\label{def:temperatura}
Para $\tau>0$, $\softmax(z,\tau)_i=e^{z_i/\tau}\big/\sum_j e^{z_j/\tau}$. En este módulo la temperatura se escribe $\tau$ para no confundirla con la longitud de contexto $T$.
\end{definicion}

\begin{proposicion}[Límites]\label{prop:limites}
Sea $z\in\R^{K}$ con un único máximo en el índice $i^{\star}$. Cuando $\tau\to 0^{+}$, $\softmax(z,\tau)\to e_{i^{\star}}$; cuando $\tau\to\infty$, $\softmax(z,\tau)\to(1/K,\dots,1/K)$.
\end{proposicion}

\begin{teorema}[La entropía crece con la temperatura]\label{teo:entropia-tau}
Sea $z\in\R^{K}$ no constante y, para $b>0$, $p^{(b)}=\softmax(bz)$, de modo que $b=1/\tau$. Entonces
\[
\frac{d}{db}\,H\bigl(p^{(b)}\bigr)=-b\,\Var_{p^{(b)}}(z)<0,
\]
donde $\Var_{p}(z)=\sum_i p_iz_i^{2}-\bigl(\sum_i p_iz_i\bigr)^{2}$. En consecuencia, $H\bigl(\softmax(z,\tau)\bigr)$ es estrictamente creciente en $\tau$.
\end{teorema}

\addtocounter{definicion}{1}%

La temperatura es la misma que la de la distribución de Boltzmann de la física estadística (módulo~1). Holtzman \emph{et al.} observan que bajar la temperatura mejora la calidad del texto generado a costa de reducir su diversidad~\cite[§3.3]{holtzman}.

\section{Muestreo top-\texorpdfstring{$k$}{k} y top-\texorpdfstring{$p$}{p}}

\begin{definicion}[Top-$k$]\label{def:topk}
Sea $1\le k\le K$ y sea $I_k(z)$ el conjunto de los índices de los $k$ mayores logits, con los empates resueltos por el índice menor. El \emph{muestreo top-$k$} usa la distribución $\softmax\bigl(\topk(z)\bigr)$, donde $\topk(z)_i=z_i$ si $i\in I_k(z)$ y $-\infty$ en otro caso.
\end{definicion}

Fan, Lewis y Dauphin introdujeron este procedimiento con $k=10$ para generar relatos y observaron que producía textos menos repetitivos que la búsqueda en haz~\cite[§5.4]{fan}. Con $k=1$ el muestreo coincide con la decodificación voraz; con $k=K$ no cambia nada. Con $z=(2;\allowbreak\,1;\allowbreak\,0{,}5;\allowbreak\,-1)$, $k=2$ y $\tau=1$, la distribución pasa a ser $(0{,}731;\allowbreak\,0{,}269;\allowbreak\,0;\allowbreak\,0)$: desaparece la cola de tokens improbables, que con 46 caracteres puede acumular una masa apreciable.

\begin{definicion}[Muestreo de núcleo o top-$p$]\label{def:topp}
Sea $q$ una distribución sobre $\V$ y $\rho\in(0,1]$. Sea $V(\rho)$ el menor conjunto de índices tal que $\sum_{j\in V(\rho)}q_j\ge\rho$. El \emph{muestreo de núcleo} usa la distribución
\[
\tilde q_i=\begin{cases} q_i\big/\sum_{j\in V(\rho)}q_j, & i\in V(\rho),\\ 0, & \text{en otro caso}.\end{cases}
\]
\end{definicion}

Es la definición de Holtzman \emph{et al.}~\cite[§3.1, ecuaciones~(2) y~(3)]{holtzman}, reproducida por CS336~\cite[§6, ecuación~(24)]{cs336}, donde el umbral se denota $p$. La distribución $q$ suele ser ya la softmax con temperatura.

\begin{proposicion}[Cálculo del núcleo]\label{prop:nucleo}
Sean $q_{(1)}\ge q_{(2)}\ge\dots\ge q_{(K)}$ las probabilidades ordenadas y $m=\min\{j:\sum_{i\le j}q_{(i)}\ge\rho\}$. Todo conjunto de masa al menos $\rho$ tiene al menos $m$ elementos, y el conjunto $S_m$ de los $m$ índices más probables alcanza esa masa. Por tanto, los conjuntos mínimos tienen exactamente $m$ elementos y $S_m$ es, entre ellos, uno de masa máxima.
\end{proposicion}

\begin{observacion}\label{obs:nucleo}
El conjunto mínimo no es único en general, aunque no haya empates: con $q=(0{,}5;\allowbreak\,0{,}3;\allowbreak\,0{,}2)$ y $\rho=0{,}6$, tanto $\{1,2\}$ como $\{1,3\}$ tienen dos elementos y masa al menos $0{,}6$. Por convenio, $V(\rho)$ designa el conjunto $S_m$, que es el que se obtiene al ordenar la distribución y acumular probabilidades hasta alcanzar $\rho$, como prescribe CS336~\cite[§6]{cs336}.
\end{observacion}

La ventaja de top-$p$ sobre top-$k$ es que el número de candidatos se adapta a la forma de la distribución. Con una distribución plana, un $k$ pequeño elimina continuaciones razonables; con una distribución muy concentrada, un $k$ grande deja entrar candidatos inapropiados, cuya probabilidad crece al renormalizar~\cite[§3.2]{holtzman}. Las interfaces de programación de modelos comerciales como Claude y Gemini exponen públicamente los parámetros de temperatura, top-$k$ y top-$p$.

\paragraph{Muestras del mini-GPT.} Con $\tau=0{,}8$ y $k=10$, casi todas las líneas generadas son correctas:
\begin{center}\small
\begin{tabular}{@{}l@{}}
\texttt{objeto. una lámpara azul de cristal. sirve para esperar sentado.}\\
\texttt{objeto. una caja azul de tela. sirve para leer de noche.}\\
\texttt{receta. hornea verde paperado.}
\end{tabular}
\end{center}
La tercera muestra el límite del modelo: mezcla un color propio de los objetos en una receta. Con $\tau=3$ la salida se degrada a símbolos sueltos, porque la distribución se acerca a la uniforme y los dígitos y signos aparecen en cualquier lugar.

\section{Efecto de los hiperparámetros}

El \cref{tab:hiper} resume entrenamientos de 800 pasos sobre el mismo corpus y con la misma semilla, en los que solo cambia lo indicado.

\begin{table}[h]
\centering\small
\caption{Efecto de los hiperparámetros en el mini-GPT (resultados del cuaderno de prácticas). Base: $d=64$, $L=2$, $n_h=4$, $T=32$, $\eta_{\max}=3\cdot 10^{-3}$. Pérdidas de validación en nats por carácter; tiempo relativo al modelo base.}\label{tab:hiper}
\begin{tabularx}{\textwidth}{@{}lrccc>{\raggedright\arraybackslash}X@{}}
\toprule
Cambio & $N$ & Paso 100 & Final & Tiempo & Lectura \\
\midrule
Base & 107\,694 & $2{,}046$ & $0{,}324$ & 1 & referencia \\
$\eta_{\max}=3\cdot 10^{-4}$ & 107\,694 & $2{,}959$ & $1{,}840$ & 1 & demasiado lenta \\
$\eta_{\max}=3\cdot 10^{-2}$ & 107\,694 & $1{,}304$ & $0{,}297$ & 1 & mejor: tolera tasas altas \\
ídem, sin calentamiento & 107\,694 & $1{,}232$ & $0{,}306$ & 1 & algo peor \\
$\eta_{\max}=3\cdot 10^{-1}$ & 107\,694 & $2{,}863$ & $2{,}244$ & 1 & inestable \\
Tasa constante & 107\,694 & $2{,}046$ & $0{,}337$ & 1 & falta el afinado final \\
$d=32$ & 29\,294 & $2{,}308$ & $0{,}649$ & $0{,}5$ & poca capacidad \\
$L=1$ & 57\,902 & $2{,}128$ & $0{,}392$ & $0{,}4$ & un bloque no basta \\
$n_h=1$ & 107\,694 & $2{,}062$ & $0{,}336$ & $0{,}7$ & algo peor \\
$T=8$ & 106\,158 & $2{,}068$ & $0{,}694$ & $0{,}3$ & no ve la etiqueta \\
$d=128$, $L=3$ & 609\,838 & $1{,}441$ & $0{,}297$ & $2{,}4$ & mejor, más caro \\
\bottomrule
\end{tabularx}
\end{table}

La tasa de aprendizaje es el factor dominante: entre $3\cdot 10^{-4}$ y $3\cdot 10^{-2}$ la pérdida final pasa de $1{,}840$ a $0{,}297$, y con $3\cdot 10^{-1}$ el entrenamiento se rompe. Goodfellow, Bengio y Courville la señalan como el hiperparámetro más importante si solo se puede ajustar uno~\cite[§11.4.1]{gbc}, y CS336 pide localizar empíricamente la tasa a partir de la cual el entrenamiento diverge y relacionarla con la mejor~\cite[§7.2, problema \emph{learning\_rate}]{cs336}. El contexto es el segundo factor: con $T=8$ la pérdida se duplica, porque el modelo pierde la información que más reduce la incertidumbre, el tipo de línea (\cref{eje:5-contexto}). Multiplicar por $5{,}7$ los parámetros da la misma mejora que subir la tasa a $3\cdot 10^{-2}$ sin coste adicional. El calentamiento y el coseno aportan poco en este modelo (unas centésimas de nat cada uno), porque es pequeño y el recorte ya lo protege; en los modelos grandes son imprescindibles.

\begin{observacion}[Reglas prácticas]\label{obs:reglas}
De lo anterior se desprenden cuatro reglas: (i) ajustar primero la tasa de aprendizaje, barriendo potencias de 10 o de $\sqrt{10}$, y quedarse con la mayor que sea estable; (ii) antes de escalar, comprobar que el modelo sobreajusta un único lote hasta pérdida casi nula; (iii) elegir $T$ de modo que cubra la dependencia que se quiere capturar, aquí una línea completa de unos 60 a 80 caracteres; (iv) más parámetros solo ayudan si hay datos y pasos suficientes para entrenarlos, como precisan las leyes de escalado del módulo~16.
\end{observacion}



\backmatter
% =====================================================================
\begin{thebibliography}{99}
\addcontentsline{toc}{chapter}{Bibliografía}
\markboth{Bibliografía}{}
\small

\bibitem{ct}
T.~M. Cover y J.~A. Thomas.
\newblock \emph{Elements of Information Theory}, 2.ª ed.
\newblock Wiley, 2006.
\newblock Cap.~2, «Entropy, Relative Entropy, and Mutual Information».

\bibitem{fan}
A.~Fan, M.~Lewis e Y.~Dauphin.
\newblock Hierarchical neural story generation.
\newblock En \emph{Proceedings of the 56th Annual Meeting of the Association for Computational Linguistics (ACL)}, 2018.
\newblock \url{https://arxiv.org/abs/1805.04833}

\bibitem{gemini}
Gemini Team, Google.
\newblock Gemini: A family of highly capable multimodal models.
\newblock Informe técnico, 2023.
\newblock \url{https://arxiv.org/abs/2312.11805}

\bibitem{gbc}
I.~Goodfellow, Y.~Bengio y A.~Courville.
\newblock \emph{Deep Learning}.
\newblock MIT Press, 2016.
\newblock Cap.~7, «Regularization for Deep Learning»; cap.~8, «Optimization for Training Deep Models»; cap.~11, «Practical Methodology».
\newblock \url{https://www.deeplearningbook.org/}

\bibitem{holtzman}
A.~Holtzman, J.~Buys, L.~Du, M.~Forbes e Y.~Choi.
\newblock The curious case of neural text degeneration.
\newblock En \emph{International Conference on Learning Representations (ICLR)}, 2020.
\newblock \url{https://arxiv.org/abs/1904.09751}

\bibitem{nanogpt}
A.~Karpathy.
\newblock nanoGPT.
\newblock Repositorio de software, GitHub, 2022--2025. Archivos \texttt{model.py}, \texttt{train.py}, \texttt{sample.py} y \texttt{config/train\_shakespeare\_char.py}. Desde noviembre de 2025 su autor lo declara obsoleto y remite a nanochat.
\newblock \url{https://github.com/karpathy/nanoGPT}

\bibitem{llama3}
Llama Team, AI @ Meta.
\newblock The Llama 3 herd of models.
\newblock Informe técnico, 2024.
\newblock \url{https://arxiv.org/abs/2407.21783}

\bibitem{sgdr}
I.~Loshchilov y F.~Hutter.
\newblock SGDR: Stochastic gradient descent with warm restarts.
\newblock En \emph{International Conference on Learning Representations (ICLR)}, 2017.
\newblock \url{https://arxiv.org/abs/1608.03983}

\bibitem{adamw}
I.~Loshchilov y F.~Hutter.
\newblock Decoupled weight decay regularization.
\newblock En \emph{International Conference on Learning Representations (ICLR)}, 2019.
\newblock \url{https://arxiv.org/abs/1711.05101}

\bibitem{gpt2}
A.~Radford, J.~Wu, R.~Child, D.~Luan, D.~Amodei e I.~Sutskever.
\newblock Language models are unsupervised multitask learners.
\newblock Informe técnico, OpenAI, 2019.

\bibitem{cs336}
Stanford CS336.
\newblock Language Modeling from Scratch.
\newblock Stanford University, primavera de 2026. Tarea~1 (\emph{basics}), «Building a Transformer LM», versión~26.0.3.
\newblock \url{https://github.com/stanford-cs336/assignment1-basics}

\bibitem{llama}
H.~Touvron \emph{et al.}
\newblock LLaMA: Open and efficient foundation language models.
\newblock Prepublicación, 2023.
\newblock \url{https://arxiv.org/abs/2302.13971}

\bibitem{vaswani}
A.~Vaswani, N.~Shazeer, N.~Parmar, J.~Uszkoreit, L.~Jones, A.~N. Gomez, Ł.~Kaiser e I.~Polosukhin.
\newblock Attention is all you need.
\newblock En \emph{Advances in Neural Information Processing Systems 30 (NIPS)}, 2017.
\newblock \url{https://arxiv.org/abs/1706.03762}

\bibitem{grok1}
xAI.
\newblock Grok-1: pesos abiertos y código de ejemplo.
\newblock Repositorio de software, GitHub, 2024.
\newblock \url{https://github.com/xai-org/grok-1}

\bibitem{cs224n}
D.~Yang y Y.~Choi.
\newblock CS224N: Natural Language Processing with Deep Learning.
\newblock Stanford University, invierno de 2026. Clase~5, «Transformers»; clase~7, «Pretraining (Scaling, Systems, Data)».
\newblock \url{https://web.stanford.edu/class/cs224n/}

\bibitem{d2l}
A.~Zhang, Z.~C. Lipton, M.~Li y A.~J. Smola.
\newblock \emph{Dive into Deep Learning}, versión 1.0.3.
\newblock Cambridge University Press, 2023.
\newblock Cap.~9, «Recurrent Neural Networks» (§9.3 y §9.5); cap.~11, «Attention Mechanisms and Transformers» (§11.9).
\newblock \url{https://d2l.ai/}

\end{thebibliography}


\end{document}
